\subsubsection{Dependencia estructural de la coordenada electrónica}

La ecuación electrónica expresa una composición jerárquica. Su primer
factor, \(\sqrt3/4\), pertenece a la estructura matricial central y fija
una normalización geométrica. El exponente \(\varphi/\pi^2\) utiliza las
coordenadas regionales ya generadas en la orientación indicada del lector
areal--volumétrico. El término \(22\alpha^3\) incorpora el acoplamiento
dodecafásico en la corrección cúbica. El factor \(1+15\Delta_4\) conserva
la normalización torsional global. Finalmente, el cociente de las secciones
de acción actúa con el exponente determinado por los lectores de retorno.
Cada operación tiene así un antecedente distinto dentro de la misma
construcción.

La articulación resulta más precisa al considerar sus dependencias.
La norma central puede calcularse antes de evaluar las coordenadas
regionales. La exponencial y la corrección cúbica necesitan esas
coordenadas y el cierre de \(\alpha\). La memoria multiplicativa necesita
además las secciones de acción y la trayectoria de \(108\) posiciones.
Por ello, conocer la expresión basal no equivale a haber evaluado la
coordenada completa. La sucesión de operaciones no agrega una masa
externa al punto \((5,5)\); desarrolla los lectores atribuidos a su
sección persistente.

La distinción entre escala y estructura proporciona también una
interpretación de la comparación física. La masa es la realización
dimensional de la coordenada completa, mientras el centro, la involución
y la vacancia describen las relaciones que la producen. Un cambio común
de unidad modifica las coordenadas dimensionales, pero deja inalteradas
las razones adimensionales y las identidades de transporte demostradas.
Éste es el sentido en que el electrón conserva su carácter estructural
al expresarse respecto de una escala cronológica.
